%% toc: P15 166 166.2
% Proofs for joint and conditional entropy, dependence, information loss, and the rate of a pair with a joint density.
\section{Joint and conditional entropy: proofs}
\label{sec:annex-p15-w02b-p15-041}

Use the definitions of Chapter~\ref{sec:p15-w02b-p15-041}. Logarithms
have a fixed base \(b>1\); \(\ln\) is the natural logarithm, so
\(\log t=\ln t/\ln b\) and \(\log e=1/\ln b\). Two observations are
used repeatedly. First, the entropy of a finite table depends only on
the list of its positive cell masses; listing the coordinates of a
tuple in another order, or adding a coordinate that is a function of
the others, changes the cells' names but not their masses. Second, for
\(u>0\),
\[
 \ln u\leq u-1,\quad\text{with equality if and only if }u=1,
 \tag{L}
\]
since \(u\mapsto u-1-\ln u\) has derivative \(1-1/u\), which is
negative on \((0,1)\) and positive on \((1,\infty)\), and the function
vanishes at \(u=1\).

\begin{annexitem}{annex:p15-w02b-p15-041-average}{thm:p15-w02b-p15-041-average}{Conditional entropy as an average}
\begin{proof}
Every positive cell has \(p(i,\cdot)>0\), and there
\(p(i,j)=p(i,\cdot)p_i(j)\). Grouping the positive cells by \(i\),
\[
 H_x(y)=-\sum_{i:\,p(i,\cdot)>0}p(i,\cdot)\sum_{j:\,p(i,j)>0}p_i(j)\log p_i(j)
       =\sum_{i:\,p(i,\cdot)>0}p(i,\cdot)H(y\mid x=i).
\]
For \(p(i,\cdot)>0\), the numbers \(p_i(j)\), \(j=1,\ldots,n\), are
nonnegative with sum \(p(i,\cdot)/p(i,\cdot)=1\). A positive
\(p_i(j)\) lies in \((0,1]\), so \(-p_i(j)\log p_i(j)\geq0\), with
equality if and only if \(p_i(j)=1\). Hence \(H(y\mid x=i)\geq0\), and
\(H(y\mid x=i)=0\) if and only if every positive \(p_i(j)\) equals one,
that is, exactly one \(j\) has \(p(i,j)>0\). The average has positive
weights \(p(i,\cdot)\), so it is nonnegative and vanishes if and only
if each of its terms vanishes.
\end{proof}
\end{annexitem}

\begin{annexitem}{annex:p15-w02b-p15-041-subadditivity}{thm:p15-w02b-p15-041-subadditivity}{Subadditivity of joint entropy}
\begin{proof}
On the finite index set of pairs \((i,j)\), put \(P(i,j)=p(i,j)\)
and \(Q(i,j)=p(i,\cdot)p(\cdot,j)\). Both are nonnegative with total
\(1\), since
\[
 \sum_{i,j}p(i,\cdot)p(\cdot,j)=\sum_ip(i,\cdot)\sum_jp(\cdot,j)=1 .
\]
On a cell with \(P(i,j)>0\) both marginals are
positive, so \(Q(i,j)>0\). By \cref{thm:p15-rmdl-gibbs},
\[
 D(P,Q)=\sum_{p(i,j)>0}p(i,j)\ln\frac{p(i,j)}{p(i,\cdot)p(\cdot,j)}\geq0,
\]
with equality if and only if \(P=Q\), which is independence. Every term
is finite, so the sum splits into three finite sums. Dividing by
\(\ln b\),
\begin{align*}
 \frac{D(P,Q)}{\ln b}
 &=\sum p(i,j)\log p(i,j)-\sum p(i,j)\log p(i,\cdot)
  -\sum p(i,j)\log p(\cdot,j)\\
 &=-H(x,y)+H(x)+H(y),
\end{align*}
by the definitions. Thus \(H(x)+H(y)-H(x,y)\geq0\), with equality if
and only if \(x\) and \(y\) are independent.
\end{proof}
\end{annexitem}

\begin{annexitem}{annex:p15-w02b-p15-041-chain}{thm:p15-w02b-p15-041-chain}{The chain identity and the conditioning bound}
\begin{proof}
On a positive cell, \(\log p_i(j)=\log p(i,j)-\log p(i,\cdot)\), and
all three numbers are finite. Multiplying by \(-p(i,j)\) and summing
over the positive cells,
\[
 H_x(y)=H(x,y)-H(x).
\]
The same computation with \(p_j(i)=p(i,j)/p(\cdot,j)\) gives
\(H_y(x)=H(x,y)-H(y)\). This proves the two identities.

\Cref{thm:p15-w02b-p15-041-average} gives \(H_x(y)\geq0\). By the
first identity and \cref{thm:p15-w02b-p15-041-subadditivity},
\[
 H_x(y)=H(x,y)-H(x)\leq H(x)+H(y)-H(x)=H(y),
\]
and equality holds exactly when \(H(x,y)=H(x)+H(y)\), that is, by the
same theorem, exactly under independence.
\end{proof}
\end{annexitem}

\begin{annexitem}{annex:p15-w02b-p15-041-averaging}{thm:p15-w02b-p15-041-averaging}{Averaging increases entropy}
\begin{proof}
Each \(p'_i\geq0\), and
\(\sum_ip'_i=\sum_jp_j\sum_ia_{ij}=\sum_jp_j=1\).

Put \(\varphi(t)=t\log t\) for \(t>0\) and \(\varphi(0)=0\), so that
\(H(p_1,\ldots,p_n)=-\sum_j\varphi(p_j)\).

\emph{Tangent inequality.} Let \(t_0>0\) and
\(\ell(t)=\varphi(t_0)+(\log t_0+\log e)(t-t_0)\). For \(t\geq0\),
\(\varphi(t)\geq\ell(t)\), with equality if and only if \(t=t_0\).
At \(t=0\), \(\varphi(0)-\ell(0)=t_0\log e>0\). For \(t>0\) put
\(u=t/t_0\); a direct expansion gives
\begin{align*}
 \varphi(t)-\ell(t)&=t\log\frac t{t_0}-(t-t_0)\log e
   =t_0\log e\,\bigl(u\ln u-u+1\bigr)\\
   &=t_0\log e\,u\Bigl(\frac1u-1-\ln\frac1u\Bigr),
\end{align*}
which is nonnegative by (L) applied to
\(1/u\) and vanishes exactly when \(u=1\).

\emph{Row inequality.} Fix \(i\). The weights \(a_{ij}\),
\(j=1,\ldots,n\), are nonnegative with sum one. If \(p'_i=0\), then
\(p_j=0\) whenever \(a_{ij}>0\), so
\(\sum_ja_{ij}\varphi(p_j)=0=\varphi(p'_i)\). If \(p'_i>0\), apply
the tangent inequality with \(t_0=p'_i\) at each \(t=p_j\), multiply
by \(a_{ij}\) and add:
\[
 \sum_ja_{ij}\varphi(p_j)\geq\varphi(p'_i)+(\log p'_i+\log e)
   \Bigl(\sum_ja_{ij}p_j-p'_i\Bigr)=\varphi(p'_i).
\]
In both cases \(\varphi(p'_i)\leq\sum_ja_{ij}\varphi(p_j)\), and
equality holds if and only if \(p_j=p'_i\) for every \(j\) with
\(a_{ij}>0\).

\emph{Summation.} Adding the row inequalities and using the column
sums,
\[
 -H(p')=\sum_i\varphi(p'_i)\leq\sum_j\varphi(p_j)\sum_ia_{ij}
   =\sum_j\varphi(p_j)=-H(p).
\]

\emph{Equality.} If \(p'\) is a rearrangement of \(p\), the two sums
of \(\varphi\) contain the same terms and the entropies agree.
Conversely, suppose \(H(p')=H(p)\). Then every row inequality is an
equality, so
\[
 a_{ij}>0\ \Longrightarrow\ p_j=p'_i .
 \tag{S}
\]
Fix a real \(c\) and put \(I_c=\{i:p'_i=c\}\) and
\(J_c=\{j:p_j=c\}\). By (S), a
positive \(a_{ij}\) with \(i\in I_c\) has \(j\in J_c\), and a positive
\(a_{ij}\) with \(j\in J_c\) has \(i\in I_c\). Using the row sums,
then these two facts, then the column sums,
\[
 |I_c|=\sum_{i\in I_c}\sum_ja_{ij}=\sum_{i\in I_c}\sum_{j\in J_c}a_{ij}
      =\sum_{j\in J_c}\sum_ia_{ij}=|J_c|.
\]
Every value therefore occurs equally often in \(p'\) and in \(p\), so
\(p'\) is a rearrangement of \(p\).

\emph{Two-coordinate transfer.} Let \(p_1<p_2\),
\(0<\delta<p_2-p_1\) and \(t=\delta/(p_2-p_1)\in(0,1)\). Take
\(a_{11}=a_{22}=1-t\), \(a_{12}=a_{21}=t\), \(a_{kk}=1\) for
\(k\geq3\), and all other entries zero. Every row and column sums to
one, \(p'_1=(1-t)p_1+tp_2=p_1+\delta\),
\(p'_2=tp_1+(1-t)p_2=p_2-\delta\), and \(p'_k=p_k\) for \(k\geq3\).
If \(p'\) were a rearrangement of \(p\), cancelling the common
coordinates \(k\geq3\) would make \(\{p_1+\delta,p_2-\delta\}\) equal
to \(\{p_1,p_2\}\) as lists with multiplicity. That is impossible,
because \(p_1<p_1+\delta<p_2\). By the equality case, the entropy
increases strictly.
\end{proof}
\end{annexitem}

\begin{annexitem}{annex:p15-w02b-p15-041-rate}{thm:p15-w02b-p15-041-rate}{Three forms of the rate and its range}
\begin{proof}
By \cref{thm:p15-w02b-p15-041-chain}, \(H_y(x)=H(x,y)-H(y)\) and
\(H_x(y)=H(x,y)-H(x)\). Substituting,
\[
 H(x)-H_y(x)=H(x)+H(y)-H(x,y)=H(y)-H_x(y).
\]
The middle expression is symmetric in \(x\) and \(y\), so
\(R(x;y)=R(y;x)\). It is nonnegative, and zero exactly under
independence, by \cref{thm:p15-w02b-p15-041-subadditivity}. Since
\(H_y(x)\geq0\) and \(H_x(y)\geq0\) by
\cref{thm:p15-w02b-p15-041-average}, the first and third forms give
\(R\leq H(x)\) and \(R\leq H(y)\). Finally \(R=H(x)\) if and only if
\(H_y(x)=0\), which by \cref{thm:p15-w02b-p15-041-average}, with the
roles of \(x\) and \(y\) exchanged, holds exactly when each value of
\(y\) of positive probability occurs in only one positive cell, that
is, when \(x\) is a function of \(y\) on the positive cells.
\end{proof}
\end{annexitem}

\begin{annexitem}{annex:p15-w02b-p15-041-correction}{thm:p15-w02b-p15-041-correction}{Conditioning on further data}
\begin{proof}
Write \(q(i,j,k)\) for the probability that \(x=i\), \(y=j\) and
\(z=k\). Apply \cref{thm:p15-w02b-p15-041-chain} three times, to the
pairs \(((x,z),y)\), \((x,(y,z))\) and \((z,y)\); by the first
observation at the head of this annex \(H((x,z),y)=H(x,y,z)\):
\begin{align*}
 H(x,y,z)&=H(y)+H_y(x,z),\\
 H(x,y,z)&=H(y,z)+H_{yz}(x),\\
 H(y,z)&=H(y)+H_y(z).
\end{align*}
Eliminating \(H(x,y,z)\) and \(H(y,z)\) gives
\(H_y(x,z)=H_y(z)+H_{yz}(x)\). Also, with \(H(x,y)=H(y)+H_y(x)\) and
the chain identity for the pair \((z,(x,y))\),
\[
 H_y(x,z)-H_y(x)=H(x,y,z)-H(x,y)=H_{xy}(z)\geq0
\]
by \cref{thm:p15-w02b-p15-041-average}. This proves the first line.
For the second, \(H_{yz}(x)=H_y(x,z)-H_y(z)\geq H_y(x)-H_y(z)\), and
\(H_y(z)\leq H(z)\) by \cref{thm:p15-w02b-p15-041-chain} for the pair
\((y,z)\).

For \(H_{yz}(x)\leq H_y(x)\), fix \(j\) with
\(r(j)\defeq\sum_{i,k}q(i,j,k)>0\) and consider the conditional table
\(q(i,j,k)/r(j)\) of \((x,z)\) given \(y=j\). Its \(x\)-marginal is
\(\sum_kq(i,j,k)/r(j)\), which is the conditional law of \(x\) given
\(y=j\); likewise its \(z\)-marginal is the conditional law of \(z\)
given \(y=j\). \Cref{thm:p15-w02b-p15-041-subadditivity} applied to
this table gives
\[
 H((x,z)\mid y=j)\leq H(x\mid y=j)+H(z\mid y=j).
\]
Multiply by \(r(j)\) and sum over such \(j\). By
\cref{thm:p15-w02b-p15-041-average}, applied with \(y\) as the
conditioning variable to the variables \((x,z)\), \(x\) and \(z\) in
turn, the result is \(H_y(x,z)\leq H_y(x)+H_y(z)\). Combined with the
identity \(H_y(x,z)=H_y(z)+H_{yz}(x)\), this gives
\(H_{yz}(x)\leq H_y(x)\).

If \(H(z)<H_y(x)\), the second line gives \(H_{yz}(x)>0\).
\end{proof}
\end{annexitem}

\begin{annexitem}{annex:p15-w02b-p15-041-loss}{thm:p15-w02b-p15-041-loss}{Coarsening and processing lose information}
\begin{proof}
(1) Put \(u=f(x)\). The table of \((x,u)\) has mass \(p(i,\cdot)\)
at \((i,f(i))\) and zero elsewhere, so it has the same positive masses
as the law of \(x\) and \(H(x,u)=H(x)\). By
\cref{thm:p15-w02b-p15-041-chain} for the pair \((u,x)\),
\[
 H(x)=H(x,u)=H(u)+H_u(x)\geq H(u),
\]
since \(H_u(x)\geq0\). Equality holds if and only if \(H_u(x)=0\). By
\cref{thm:p15-w02b-p15-041-average}, with \(u\) as conditioning
variable, this happens exactly when each value \(k\) of \(u\) of
positive probability has only one \(i\) with \(p(i,\cdot)>0\) and
\(f(i)=k\). That is the statement that \(f\) is one-to-one on
\(\{i:p(i,\cdot)>0\}\).

(2) Put \(u=g(y)\). The table of \((x,u,y)\) has mass \(p(i,j)\) at
\((i,g(j),j)\) and zero elsewhere, and the table of \((u,y)\) has mass
\(p(\cdot,j)\) at \((g(j),j)\). These have the same positive masses as
the tables of \((x,y)\) and \(y\), so by the chain identity
\[
 H_{uy}(x)=H(x,u,y)-H(u,y)=H(x,y)-H(y)=H_y(x).
\]
The last inequality of \cref{thm:p15-w02b-p15-041-correction}, with
\(u\) in the place of \(y\) and \(y\) in the place of \(z\), gives
\(H_{uy}(x)\leq H_u(x)\). Hence \(H_y(x)\leq H_u(x)\) and
\[
 R(x;u)=H(x)-H_u(x)\leq H(x)-H_y(x)=R(x;y).\qedhere
\]
\end{proof}
\end{annexitem}

\begin{annexitem}{annex:p15-w02b-p15-041-density-rate}{thm:p15-w02b-p15-041-density-rate}{The rate of a pair with a joint density}
\begin{proof}
Since \(H(x)\), \(H(u)\) and \(H(x,u)\) are finite,
\cref{thm:p15-w02b-p15-042-conditional} shows that \(H_x(u)\) and
\(H_u(x)\) are finite and that
\[
 H(x,u)=H(x)+H_x(u)=H(u)+H_u(x).
\]
The second equality gives \(H_u(x)=H(x,u)-H(u)\), hence
\(H(x)-H_u(x)=H(x)+H(u)-H(x,u)\). The first gives
\(H_x(u)=H(x,u)-H(x)\), hence \(H(u)-H_x(u)=H(x)+H(u)-H(x,u)\). All
three expressions are equal. The last is unchanged when the roles of
\(x\) and \(u\) are exchanged, and \(R(u;x)\) is by definition
\(H(u)-H_x(u)\); so \(R(x;u)=R(u;x)\).

Because \(H(x)\) and \(H(u)\) are finite,
\cref{thm:p15-w02b-p15-042-joint} gives \(H(x,u)\leq H(x)+H(u)\), with
equality if and only if \(p(x,u)=p(x)p(u)\) for almost every \((x,u)\).
By the third form, \(R(x;u)\geq0\), with equality exactly in that case,
which is independence in the sense of
\cref{def:p15-w02b-p15-041-density}.
\end{proof}
\end{annexitem}

\begin{annexitem}{annex:p15-w02b-p15-041-coordinates}{thm:p15-w02b-p15-041-coordinates}{Recoding one variable of a pair}
\begin{proof}
\emph{Item 1.} The hypotheses of
\cref{thm:p15-w02b-p15-042-coordinate-change-thm} hold: \((T,j)\) is a
coordinate change, \(P(x\in U)=1\), \(H(x)\) is defined because it is
finite, and \(\int p\,|\log j|\,dx<\infty\). That theorem gives a
density for \(y\), namely \(p(T^{-1}t)/j(T^{-1}t)\) on \(V\) and zero
elsewhere, and
\[
 H(y)=H(x)+\int p(x)\log j(x)\,dx .
\]
The integral is finite because its absolute value is at most
\(\int p\,|\log j|\,dx\), and \(H(x)\) is finite, so \(H(y)\) is
finite. Since \(|J(x/y)|=1/j(x)\), \(\log j(x)=-\log|J(x/y)|\), which
gives the second form.

\emph{Item 2: the pair has a coordinate change.} The identity map of
\(\mathbb R^k\) is a coordinate change with factor \(1\), by item~1 of
\cref{def:p15-w02b-p15-042-coordinate-change} with the identity matrix
and \(b=0\). By item~3 of that definition, which holds by
\cref{thm:p15-w02b-p15-042-linear-lebesgue}, the map
\(S(x,u)=(T(x),u)\) from \(U\times\mathbb R^k\) onto
\(V\times\mathbb R^k\), with factor \(j(x)\cdot1=j(x)\), is a
coordinate change in \(\mathbb R^{n+k}\). Also
\(P((x,u)\in U\times\mathbb R^k)=P(x\in U)=1\).

\emph{Item 2: integrability.} The marginal density of \(x\) is
\(p(x)=\int p(x,u)\,du\). By Tonelli's theorem,
\[
 \iint p(x,u)\,|\log j(x)|\,dx\,du=\int|\log j(x)|\,p(x)\,dx<\infty ,
\]
and then, by Fubini's theorem,
\(\iint p(x,u)\log j(x)\,dx\,du=\int p(x)\log j(x)\,dx=c\).

\emph{Item 2: the joint entropy.} Apply
\cref{thm:p15-w02b-p15-042-coordinate-change-thm} to the vector
\((x,u)\) in \(\mathbb R^{n+k}\), the coordinate change \(S\), and the
finite entropy \(H(x,u)\). The pair \((y,u)=S(x,u)\) has the joint
density \(q(t,u)=p(T^{-1}t,u)/j(T^{-1}t)\) for \(t\in V\), zero
elsewhere, and
\[
 H(y,u)=H(x,u)+\iint p(x,u)\log j(x)\,dx\,du=H(x,u)+c ,
\]
which is finite.

\emph{Item 2: the marginals.} Integrating \(q(t,u)\) over \(u\) gives,
for \(t\in V\), \(p(T^{-1}t)/j(T^{-1}t)\), the density of \(y\) found in
item~1. The second variable of the pair is still \(u\), whose density is
\(p(u)\) and whose entropy \(H(u)\) is finite. So \((y,u)\) satisfies
\cref{def:p15-w02b-p15-041-density} with \(H(y)\), \(H(u)\) and
\(H(y,u)\) finite.

\emph{Item 2: the conditional entropies and the rate.} By
\cref{thm:p15-w02b-p15-042-conditional}, applied to both pairs,
\[
 H_u(y)=H(y,u)-H(u)=H(x,u)+c-H(u)=H_u(x)+c,
\]
\[
 H_y(u)=H(y,u)-H(y)=\bigl(H(x,u)+c\bigr)-\bigl(H(x)+c\bigr)=H_x(u).
\]
Finally
\(R(y;u)=H(y)-H_u(y)=\bigl(H(x)+c\bigr)-\bigl(H_u(x)+c\bigr)=R(x;u)\).
\end{proof}
\end{annexitem}
